\section{Reactor Plant Physics and Modeling Details}
\label{ap:plantmodel}

This appendix documents the reactor plant model used in the present study.  The implemented Modelica system is a point-kinetics reactor model coupled to a four-node lumped thermal network representing the fuel, moderator/graphite, heat-pipe, and nitrogen-loop temperatures.  In the current implementation, the secondary side is not modeled as an additional dynamic thermal node; instead, steam-side heat removal is represented through an algebraic heat-transfer relation at a prescribed steam outlet temperature.  The control input is a common drum-angle signal applied uniformly to all drums.

\subsection{Governing equations}

The model state consists of normalized reactor power $n$, six delayed-neutron precursor concentrations $c_i$, and four thermal states $T_f$, $T_m$, $T_{hp}$, and $T_{N}$.  Thermal power is obtained from normalized power through
\begin{equation}
    P(t) = n(t)\,P_r,
    \qquad
    P_{\mathrm{MW}}(t) = \frac{P(t)}{10^6}.
\end{equation}
The total reactivity is decomposed into drum, fuel-temperature, and moderator-temperature contributions:
\begin{align}
    \rho_{\mathrm{drums}}(t)
    &=
    n_d \rho_{\max}^{(1)}
    \frac{1-\cos\!\left(\frac{\pi}{180}\theta_{\mathrm{drum}}(t)\right)}{2}
    - \rho_{\mathrm{ss,tot}},
    \\
    \rho_f(t)
    &= \alpha_f \bigl(T_f(t)-T_{f0}\bigr),
    \\
    \rho_m(t)
    &= \alpha_m \bigl(T_m(t)-T_{m0}\bigr),
    \\
    \rho(t)
    &= \rho_{\mathrm{drums}}(t) + \rho_f(t) + \rho_m(t).
\end{align}
Reactivity in dollars is defined as
\begin{equation}
    \rho_{\$}(t) = \frac{\rho(t)}{\beta}.
\end{equation}

The neutronics are governed by the standard six-group point-kinetics equations,
\begin{align}
    \frac{dn}{dt}
    &=
    \left(\frac{\rho(t)-\beta}{\Lambda}\right)n(t)
    +
    \sum_{i=1}^{6} \lambda_i c_i(t),
    \\
    \frac{dc_i}{dt}
    &=
    \frac{\beta_i}{\Lambda}n(t) - \lambda_i c_i(t),
    \qquad i=1,\dots,6.
\end{align}

The thermal network is written as a set of lumped energy balances.  Fuel temperature evolves according to
\begin{equation}
    \frac{dT_f}{dt}
    =
    \frac{f_h\,P(t)}{M_f c_{p,f}}
    +
    \frac{T_m(t)-T_f(t)}{\tau_{f\to g}},
\end{equation}
while the moderator/graphite node satisfies
\begin{equation}
    \frac{dT_m}{dt}
    =
    \frac{(1-f_h)\,P(t)}{M_g c_{p,g}}
    +
    \frac{T_f(t)-T_m(t)}{\tau_{g\to f}}
    +
    \frac{T_{hp}(t)-T_m(t)}{\tau_{g\to hp}}.
\end{equation}
The heat-pipe node evolves as
\begin{equation}
    \frac{dT_{hp}}{dt}
    =
    \frac{T_m(t)-T_{hp}(t)}{\tau_{hp\to g}}
    +
    \frac{T_{N}(t)-T_{hp}(t)}{\tau_{hp\to N}}.
\end{equation}
Heat rejected to the steam side is modeled algebraically by
\begin{align}
    Q_{\mathrm{to\,steam}}(t)
    &= G_{N\text{-}sg}\bigl(T_{N}(t)-T_{\mathrm{steam,out}}\bigr),
    \\
    \dot{m}_{\mathrm{steam}}(t)
    &= \frac{Q_{\mathrm{to\,steam}}(t)}{\Delta h},
\end{align}
and the nitrogen-loop temperature obeys
\begin{equation}
    \frac{dT_{N}}{dt}
    =
    \frac{T_{hp}(t)-T_{N}(t)}{\tau_{N\to hp}}
    -
    \frac{Q_{\mathrm{to\,steam}}(t)}{M_{N}c_{p,N}}.
\end{equation}

The initial conditions are chosen as the nominal steady-state operating point,
\begin{equation}
    n(0)=1,
    \qquad
    c_i(0)=\frac{\beta_i}{\Lambda\lambda_i},
    \qquad
    T_f(0)=T_{f0},
    \qquad
    T_m(0)=T_{m0},
    \qquad
    T_{hp}(0)=T_{hp0},
    \qquad
    T_{N}(0)=T_{N2,0}.
\end{equation}

\subsection{Neutronics and control parameters}

\begin{table}[htbp]
\centering
\caption{Neutronics and control parameters used in \texttt{HPMicroPK}.}
\label{tab:plant-neutronics}
\begin{tabular}{llll}
\toprule
Group & Symbol & Value & Unit \\
\midrule
Prompt generation & $\Lambda$ & $1.95734\times 10^{-4}$ & s \\
Delayed groups & $N_g$ & $6$ & -- \\
Delayed decay constants &
$\lambda_i$ &
$\{0.01334,\;0.03274,\;0.1208,\;0.3028,\;0.8495,\;2.853\}$ &
s$^{-1}$ \\
Delayed fractions &
$\beta_i$ &
$\{2.90591,\;11.2869,\;11.6402,\;26.9454,\;9.20043,\;4.41185\}\times 10^{-4}$ &
-- \\
Total delayed fraction & $\beta$ & $6.639069\times 10^{-3}$ & -- \\
Steady precursor concentrations &
$C_{i,\mathrm{ss}}$ &
$\{111.2910,\;176.1285,\;49.2297,\;45.4635,\;5.53323,\;0.790047\}$ &
-- \\
Rated thermal power & $P_r$ & $6.0\times 10^{6}$ & W \\
Fuel power fraction & $f_h$ & $0.90$ & -- \\
Fuel feedback & $\alpha_f$ & $-4.59\times 10^{-5}$ & $\Delta k/k$/K \\
Moderator feedback & $\alpha_m$ & $-2.21\times 10^{-5}$ & $\Delta k/k$/K \\
Number of drums & $n_d$ & $12$ & -- \\
Total drum worth & $\rho_{\max,\mathrm{tot}}$ & $-13174$ & pcm \\
Per-drum worth & $\rho_{\max}^{(1)}$ & $-1.0978333\times 10^{-2}$ & $\Delta k/k$ \\
Reference drum angle & $u_0$ & $45.0$ & deg \\
Single-drum steady reactivity & $\rho_{\mathrm{ss}}^{(1)}$ & $-1.6077397\times 10^{-3}$ & $\Delta k/k$ \\
Total steady drum reactivity & $\rho_{\mathrm{ss,tot}}$ & $-1.9292876\times 10^{-2}$ & $\Delta k/k$ \\
\bottomrule
\end{tabular}
\end{table}

\subsection{Thermal and physical parameters}

\begin{longtable}{llll}
\caption{Thermal design values, lumped properties, and derived parameters used in the four-node thermal network.}
\label{tab:plant-thermal}
\\
\toprule
Group & Symbol & Value & Unit \\
\midrule
\endfirsthead

\toprule
Group & Symbol & Value & Unit \\
\midrule
\endhead

\bottomrule
\endfoot

\multicolumn{4}{l}{\textbf{Reference temperatures}} \\
Fuel & $T_{f,0}$ & $1173.15$ & K \\
Moderator/graphite & $T_{m,0}$ & $1150.15$ & K \\
Heat pipe & $T_{hp,0}$ & $1073.15$ & K \\
Nitrogen inlet & $T_{N,\mathrm{in}}$ & $683.15$ & K \\
Nitrogen outlet & $T_{N,\mathrm{out}}$ & $1073.15$ & K \\
Nitrogen reference average & $T_{N,0}$ & $878.15$ & K \\
Steam-generator boundary & $T_{sg,0}$ & $505.15$ & K \\
Feedwater inlet & $T_{fw,\mathrm{in}}$ & $471.45$ & K \\
Steam outlet boundary & $T_{\mathrm{steam,out}}$ & $505.15$ & K \\

\addlinespace
\multicolumn{4}{l}{\textbf{Fuel lump}} \\
Fuel mass & $M_f$ & $4355.259$ & kg \\
Fuel specific heat & $c_{p,f}$ & $748.72$ & J/(kg$\cdot$K) \\
Fuel inverse heat capacity & $(M_f c_{p,f})^{-1}$ & $3.06768\times 10^{-7}$ & K/J \\

\addlinespace
\multicolumn{4}{l}{\textbf{Moderator/graphite lump}} \\
Moderator mass & $M_g$ & $673.921$ & kg \\
Moderator specific heat & $c_{p,g}$ & $1850.0$ & J/(kg$\cdot$K) \\

\addlinespace
\multicolumn{4}{l}{\textbf{Heat-pipe lump}} \\
Heat-pipe mass & $M_{hp}$ & $119.678$ & kg \\
Heat-pipe specific heat & $c_{p,hp}$ & $771.3$ & J/(kg$\cdot$K) \\

\addlinespace
\multicolumn{4}{l}{\textbf{Nitrogen-loop lump}} \\
Nitrogen specific heat & $c_{p,N}$ & $1182.0$ & J/(kg$\cdot$K) \\
Nitrogen residence time & $\tau_N$ & $5.0$ & s \\
Nitrogen mass flow & $\dot{m}_N$ & $13.015749$ & kg/s \\
Nitrogen lumped mass & $M_{N}$ & $65.078745$ & kg \\

\addlinespace
\multicolumn{4}{l}{\textbf{Steam-side effective quantity}} \\
Effective enthalpy rise & $\Delta h$ & $2.03587\times 10^{6}$ & J/kg \\

\addlinespace
\multicolumn{4}{l}{\textbf{Derived conductances}} \\
Fuel--moderator & $G_{f\text{-}g}$ & $2.3478261\times 10^{5}$ & W/K \\
Moderator--heat pipe & $G_{g\text{-}hp}$ & $7.7922078\times 10^{4}$ & W/K \\
Heat pipe--nitrogen & $G_{hp\text{-}N}$ & $3.0769231\times 10^{4}$ & W/K \\
Nitrogen--steam boundary & $G_{N\text{-}sg}$ & $1.6085791\times 10^{4}$ & W/K \\

\addlinespace
\multicolumn{4}{l}{\textbf{Model time constants}} \\
Fuel response to moderator & $\tau_{f\to g}$ & $13.888889$ & s \\
Moderator response to fuel & $\tau_{g\to f}$ & $5.310248$ & s \\
Moderator response to heat pipe & $\tau_{g\to hp}$ & $16.000008$ & s \\
Heat-pipe response to moderator & $\tau_{hp\to g}$ & $1.184615$ & s \\
Heat-pipe response to nitrogen & $\tau_{hp\to N}$ & $2.999998$ & s \\
Nitrogen response to heat pipe & $\tau_{N\to hp}$ & $2.500000$ & s \\
Nitrogen response to steam boundary & $\tau_{N\to sg}$ & $4.782051$ & s \\

\end{longtable}